\scanpage{116}
Considérons l'algèbre produit tensoriel
\[
 A_1\otimes A_2\otimes A_3;
\]
\unsure{Via} les $\sigma_{A_i}$, le groupe $\mathbb F_2^3$ opère dessus. Il n'y a pas en général d'opération de ce groupe sur $A$, compatible avec
\[
 \bigotimes A_i\mathbin{-}A,
\]
\deleted{car} il faudrait qu'il y ait p.ex. un autom. $\tau_1$ de $A$ qui induise l'identité sur $A_2,A_3$, et $\sigma_{A_1}$ sur $A_1$; Mais dans les cas les plus simples, $A$ est engendré par $A_2,A_3$, donc $\tau_1$ devrait être l'identité, alors qu'il induit $\sigma_{A_1}$ sur $A_1$, qui ne l'est pas.

Par contre, \deleted{soit}
\[
 \begin{aligned}
 \unsure{\gamma|\Gamma}\deleted{\unsure{=}}&\subset\mathbb F_2^3\\
 &\mathrel{\text{''}}\operatorname{Ker}(\mathbb F_2^3\longrightarrow\mathbb F_2)\\
 &(\xi_i)\mathbin{-}\textstyle\sum\xi_i
 \end{aligned}
\]
formé des éléments $(\xi_1,\xi_2,\xi_3)$ tels que le nb des composantes $\ne0$ (i.e. $=1$) soit $0$ ou $2$ i.e. soit pair. \deleted{\unsure{Donc}} il y \deleted{\unsure{en}} a trois éléments non nuls de ce groupe, savoir
\[
 \deleted{(0,0,1)},\qquad\begin{gathered}(0,1,1)\\(1,0,1)\\(1,1,0)\end{gathered}
\]
qui correspondent aux autom
\[
 \begin{aligned}
 \mathrm{id}_{A_1}\otimes\sigma_{A_2}\otimes\sigma_{A_3}&\overset{\mathrm{déf}}{=}\tau_1\\
 \sigma_{A_1}\otimes\mathrm{id}_{A_2}\otimes\sigma_{A_3}&\overset{\mathrm{déf}}{=}\tau_2\\
 \sigma_{A_1}\otimes\sigma_{A_2}\otimes\mathrm{id}_{A_3}&\overset{\mathrm{déf}}{=}\tau_3
 \end{aligned}
\]
\mspage{115}
